\documentclass[11pt]{article}
\usepackage[a4paper,margin=1in]{geometry}
\usepackage[T1]{fontenc}
\usepackage[utf8]{inputenc}
\usepackage{lmodern}
\usepackage{amsmath,amssymb}
\usepackage{enumitem}
\usepackage{hyperref}
\usepackage{xcolor}

\hypersetup{
  colorlinks=true,
  linkcolor=blue,
  urlcolor=blue
}

\title{Latto-Latto Control Research Changelog}
\author{}
\date{Last updated: 2026-05-22}

\begin{document}

\maketitle

This document records major changes to the Latto-latto model, reward design, and training setup.

\paragraph{Rule for future updates.}
Any major change to the model, reward, controller, training setup, or evaluation protocol should be added to this file.

\section*{2026-05-22}

\subsection*{Change 1: Reward Updated to Penalize Vertical Drift}

\paragraph{Motivation.}
\begin{itemize}[leftmargin=1.5em]
  \item The original sparse reward was able to produce bouncing behavior, but the pivot vertical position $z$ could drift over time.
  \item This made the learned behavior less physically meaningful and less suitable for fair controller comparison.
\end{itemize}

\paragraph{Previous reward.}
\[
r_t =
\begin{cases}
1, & \text{if a collision occurs and the collision side alternates}, \\
0, & \text{otherwise.}
\end{cases}
\]

\paragraph{Updated reward.}
\[
r_t = r_t^{\mathrm{collision}} - \alpha z_t^2
\]

with

\[
r_t^{\mathrm{collision}} =
\begin{cases}
1, & \text{if a collision occurs and the collision side alternates}, \\
0, & \text{otherwise.}
\end{cases}
\]

where:
\begin{itemize}[leftmargin=1.5em]
  \item $z_t$ is the vertical position of the pivot.
  \item $\alpha$ is the vertical position penalty weight.
\end{itemize}

\paragraph{Current setting in code.}
\[
\alpha = 0.1
\]

\paragraph{Implementation notes.}
\begin{itemize}[leftmargin=1.5em]
  \item The reward is implemented in \texttt{latto\_latto\_model.py}.
  \item The environment now exposes reward components through the \texttt{info} dictionary:
  \begin{itemize}[leftmargin=1.5em]
    \item \texttt{sparse\_reward}
    \item \texttt{z\_position\_penalty}
  \end{itemize}
\end{itemize}

\paragraph{Observed outcome.}
\begin{itemize}[leftmargin=1.5em]
  \item PPO training completed successfully with the modified reward.
  \item The learned policy strongly reduced vertical drift.
  \item In quick evaluation, the policy also collapsed to a near-stationary solution with zero alternating collisions.
\end{itemize}

\paragraph{Interpretation.}
\begin{itemize}[leftmargin=1.5em]
  \item The drift penalty successfully regularized the base motion.
  \item The penalty weight $\alpha = 0.1$ appears too strong relative to the sparse collision reward.
  \item A smaller penalty weight, such as $0.01$ or $0.02$, should be tested next.
\end{itemize}

\subsection*{Change 2: Separate PPO Checkpoint for Reward-Penalty Experiment}

\paragraph{Motivation.}
Preserve the original PPO checkpoint while training a new policy with the modified reward.

\paragraph{Change.}
Training output path changed from \texttt{ppo\_latto} to \texttt{ppo\_latto\_z\_penalty}.

\paragraph{Result.}
The reward-penalty experiment can now be evaluated independently from the original sparse-reward PPO model.

\subsection*{Change 3: Testing Script Updated to Load the New Checkpoint}

\paragraph{Motivation.}
Keep the default evaluation script aligned with the most recent model variant.

\paragraph{Change.}
\texttt{testing\_latto.py} now loads \texttt{ppo\_latto\_z\_penalty}.

\paragraph{Note.}
If multiple controller or reward variants are added later, this script should eventually accept a model path as an argument instead of hardcoding one checkpoint.

\subsection*{Change 4: Visualization Workflow Improved}

\paragraph{Motivation.}
The rendering logic was previously tied to repeated global \texttt{matplotlib} calls, which was less stable for repeated runs.

\paragraph{Change.}
\begin{itemize}[leftmargin=1.5em]
  \item Rendering now uses persistent \texttt{matplotlib} figure and axes handles.
  \item The environment now includes a \texttt{close()} method to release plotting resources cleanly.
  \item \texttt{testing\_latto.py} now handles interactive and non-interactive backends more safely.
\end{itemize}

\paragraph{Research relevance.}
This is not a dynamics change, but it improves repeatability of visual inspection during controller evaluation.

\subsection*{Change 5: Testing Script Now Plots Pose Trajectories}

\paragraph{Motivation.}
\begin{itemize}[leftmargin=1.5em]
  \item The previous testing workflow only saved the control input history.
  \item For controller comparison, it is also important to inspect how the pivot position and pendulum angle evolve over time.
\end{itemize}

\paragraph{Change.}
\begin{itemize}[leftmargin=1.5em]
  \item \texttt{testing\_latto.py} now records and plots:
  \begin{itemize}[leftmargin=1.5em]
    \item pivot vertical position $z(t)$
    \item pendulum angle $\theta(t)$
  \end{itemize}
  \item The script now saves an additional figure, \texttt{pose.png}.
\end{itemize}

\paragraph{Research relevance.}
This improves the evaluation workflow by making drift, oscillation amplitude, and near-stationary collapse easier to detect visually.

\subsection*{Change 6: Parameterized Alpha Sweep for Reward Comparison}

\paragraph{Motivation.}
\begin{itemize}[leftmargin=1.5em]
  \item The first drift-penalty experiment with $\alpha = 0.1$ improved vertical centering, but it suppressed alternating collisions too strongly.
  \item A fair next step is to compare several smaller $\alpha$ values under the same PPO setup.
\end{itemize}

\paragraph{Change.}
\begin{itemize}[leftmargin=1.5em]
  \item \texttt{latto\_latto\_model.py} now accepts the vertical position penalty weight as a constructor argument.
  \item Added \texttt{alpha\_sweep\_experiment.py} to train and evaluate PPO models for:
  \begin{itemize}[leftmargin=1.5em]
    \item $\alpha = 0.01$
    \item $\alpha = 0.02$
    \item $\alpha = 0.05$
  \end{itemize}
  \item The sweep script saves:
  \begin{itemize}[leftmargin=1.5em]
    \item per-alpha PPO checkpoints
    \item \texttt{alpha\_sweep\_results.json}
    \item \texttt{alpha\_pose\_comparison.png}
  \end{itemize}
\end{itemize}

\paragraph{Research relevance.}
This turns the reward-tuning question into a reproducible experiment and directly supports comparison of $z(t)$ and $\theta(t)$ across reward weights.

\subsection*{Change 7: Collision Dynamics Updated from Ideal Reflection to Inelastic Impact}

\paragraph{Motivation.}
\begin{itemize}[leftmargin=1.5em]
  \item The previous collision model used a perfect sign flip of angular velocity.
  \item That assumption preserved too much energy at impact and made the hybrid dynamics less realistic.
  \item For a stronger paper, the collision event should include explicit impact loss.
\end{itemize}

\paragraph{Previous collision reset.}
\[
\dot{\theta}^{+} = - \dot{\theta}^{-}
\]

\paragraph{Updated collision reset.}
\[
\dot{\theta}^{+} = - e \dot{\theta}^{-}
\]

where:
\begin{itemize}[leftmargin=1.5em]
  \item $\dot{\theta}^{-}$ is the angular velocity immediately before impact.
  \item $\dot{\theta}^{+}$ is the angular velocity immediately after impact.
  \item $e \in (0, 1]$ is the coefficient of restitution.
\end{itemize}

\paragraph{Current setting in code.}
\[
e = 0.9
\]

\paragraph{Impact energy loss proxy.}
\[
\Delta E_{\mathrm{impact}}
=
\frac{1}{2}\left((\dot{\theta}^{-})^{2} - (\dot{\theta}^{+})^{2}\right)
\]

\paragraph{Implementation notes.}
\begin{itemize}[leftmargin=1.5em]
  \item \texttt{latto\_latto\_model.py} now accepts \texttt{collision\_restitution} as a constructor argument.
  \item The environment now reports additional collision diagnostics through \texttt{info}:
  \begin{itemize}[leftmargin=1.5em]
    \item \texttt{collision\_flag}
    \item \texttt{pre\_collision\_theta\_dot}
    \item \texttt{post\_collision\_theta\_dot}
    \item \texttt{collision\_energy\_loss}
  \end{itemize}
\end{itemize}

\paragraph{Research relevance.}
This makes the environment a more meaningful hybrid impact-control benchmark. It also opens a stronger next experiment: compare controller robustness across different restitution values.

\subsection*{Change 8: Alpha Sweep Repeated on the Inelastic-Collision Model}

\paragraph{Motivation.}
\begin{itemize}[leftmargin=1.5em]
  \item After introducing inelastic impacts with restitution $e = 0.9$, the previous alpha-sweep results no longer reflected the current model.
  \item The reward comparison needed to be repeated on the updated hybrid dynamics.
\end{itemize}

\paragraph{Change.}
\begin{itemize}[leftmargin=1.5em]
  \item \texttt{alpha\_sweep\_experiment.py} was updated to run the sweep on the current model with:
  \begin{itemize}[leftmargin=1.5em]
    \item $e = 0.9$
    \item $\alpha = 0.01$
    \item $\alpha = 0.02$
    \item $\alpha = 0.05$
  \end{itemize}
  \item Separate artifacts are now saved for the inelastic sweep:
  \begin{itemize}[leftmargin=1.5em]
    \item \texttt{ppo\_latto\_inelastic\_alpha\_0p010.zip}
    \item \texttt{ppo\_latto\_inelastic\_alpha\_0p020.zip}
    \item \texttt{ppo\_latto\_inelastic\_alpha\_0p050.zip}
    \item \texttt{alpha\_sweep\_inelastic\_results.json}
    \item \texttt{alpha\_pose\_inelastic\_comparison.png}
  \end{itemize}
\end{itemize}

\paragraph{Observed outcome.}
\begin{itemize}[leftmargin=1.5em]
  \item $\alpha = 0.01$: \texttt{max |z| = 0.0250}, \texttt{mean |z| = 0.0092}, alternating collisions = \texttt{0}
  \item $\alpha = 0.02$: \texttt{max |z| = 0.0409}, \texttt{mean |z| = 0.0256}, alternating collisions = \texttt{0}
  \item $\alpha = 0.05$: \texttt{max |z| = 0.0313}, \texttt{mean |z| = 0.0121}, alternating collisions = \texttt{0}
\end{itemize}

\paragraph{Interpretation.}
\begin{itemize}[leftmargin=1.5em]
  \item $\alpha = 0.01$ remained the best of the three on vertical centering for this rollout.
  \item All three policies still collapsed to non-bouncing behavior in deterministic evaluation.
  \item The inelastic collision model made the task more realistic, but it did not by itself recover the desired alternating-impact rhythm under the current reward design.
\end{itemize}

\paragraph{Research relevance.}
This is a useful negative result. It suggests that model fidelity alone is not enough; the reward or controller structure likely also needs to change to sustain rhythmic impacts under dissipative collisions.

\subsection*{Change 9: Dead-Zone Reward for Vertical Motion and \texttt{z\_tol} Sweep}

\paragraph{Motivation.}
\begin{itemize}[leftmargin=1.5em]
  \item Penalizing all vertical motion in $z$ reduced the ability of the controller to pump energy into the swing.
  \item A better reward should allow $z$ to move freely within a small working band and only penalize excessive drift outside that band.
\end{itemize}

\paragraph{Previous vertical penalty.}
\[
r_t = r_t^{\mathrm{collision}} - \alpha z_t^2
\]

\paragraph{Updated vertical penalty with dead zone.}
\[
r_t = r_t^{\mathrm{collision}} - \alpha \max(0, |z_t| - z_{\mathrm{tol}})^2
\]

where:
\begin{itemize}[leftmargin=1.5em]
  \item $z_{\mathrm{tol}}$ is the allowed free-motion range for the pivot.
  \item No $z$ penalty is applied while $|z_t| \le z_{\mathrm{tol}}$.
\end{itemize}

\paragraph{Current sweep setting.}
\begin{itemize}[leftmargin=1.5em]
  \item $\alpha = 0.1$
  \item $e = 0.9$
  \item $z_{\mathrm{tol}} \in \{0.03, 0.05, 0.08\}$
\end{itemize}

\paragraph{Implementation notes.}
\begin{itemize}[leftmargin=1.5em]
  \item \texttt{latto\_latto\_model.py} now accepts \texttt{z\_position\_tolerance}.
  \item The environment now reports \texttt{z\_position\_excess} through \texttt{info}.
  \item Added \texttt{ztol\_sweep\_experiment.py} to retrain and compare PPO across tolerance values.
\end{itemize}

\paragraph{Research relevance.}
This is a more physically meaningful regularization scheme because it distinguishes between useful excitation motion and undesirable long-term drift.

\paragraph{Observed outcome.}
\begin{itemize}[leftmargin=1.5em]
  \item $z_{\mathrm{tol}} = 0.03$: \texttt{max |z| = 0.0302}, \texttt{mean |z| = 0.0160}, alternating collisions = \texttt{0}
  \item $z_{\mathrm{tol}} = 0.05$: \texttt{max |z| = 0.1025}, \texttt{mean |z| = 0.0773}, alternating collisions = \texttt{0}
  \item $z_{\mathrm{tol}} = 0.08$: \texttt{max |z| = 0.0469}, \texttt{mean |z| = 0.0191}, alternating collisions = \texttt{0}
\end{itemize}

\paragraph{Interpretation.}
\begin{itemize}[leftmargin=1.5em]
  \item The dead-zone reward clearly changed the learned $z$ behavior compared with the everywhere-penalized $z^2$ reward.
  \item $z_{\mathrm{tol}} = 0.03$ produced the tightest vertical bound, while $z_{\mathrm{tol}} = 0.08$ allowed moderate motion but still stayed within the free band.
  \item $z_{\mathrm{tol}} = 0.05$ was the weakest of the three on vertical containment in this evaluation.
  \item None of the tested tolerance values recovered alternating-collision behavior under the current reward structure.
\end{itemize}

\subsection*{Change 10: Sweep of Collision Reward Weight $w_c$}

\paragraph{Motivation.}
\begin{itemize}[leftmargin=1.5em]
  \item The dead-zone reward improved how $z$ was treated, but alternating-collision behavior still did not reappear.
  \item This suggested that the collision reward itself might be too weak relative to the easier low-motion strategy.
\end{itemize}

\paragraph{Updated reward.}
\[
r_t =
w_c \, r_t^{\mathrm{collision}}
- \alpha \max(0, |z_t| - z_{\mathrm{tol}})^2
\]

where:
\begin{itemize}[leftmargin=1.5em]
  \item $w_c$ is the collision reward weight.
  \item $r_t^{\mathrm{collision}} = 1$ for a valid alternating collision and $0$ otherwise.
\end{itemize}

\paragraph{Current sweep setting.}
\begin{itemize}[leftmargin=1.5em]
  \item $\alpha = 0.1$
  \item $z_{\mathrm{tol}} = 0.08$
  \item $e = 0.9$
  \item $w_c \in \{1, 2, 5\}$
\end{itemize}

\paragraph{Implementation notes.}
\begin{itemize}[leftmargin=1.5em]
  \item \texttt{latto\_latto\_model.py} now accepts \texttt{collision\_reward\_weight}.
  \item Added \texttt{collision\_reward\_sweep\_experiment.py} to retrain and compare PPO across collision-reward weights.
\end{itemize}

\paragraph{Research relevance.}
This isolates whether the main issue is insufficient task incentive rather than only over-regularization of pivot motion.

\paragraph{Observed outcome.}
\begin{itemize}[leftmargin=1.5em]
  \item $w_c = 1$: \texttt{max |z| = 0.0318}, \texttt{mean |z| = 0.0160}, alternating collisions = \texttt{0}
  \item $w_c = 2$: \texttt{max |z| = 0.0275}, \texttt{mean |z| = 0.0161}, alternating collisions = \texttt{0}
  \item $w_c = 5$: \texttt{max |z| = 0.0070}, \texttt{mean |z| = 0.0035}, alternating collisions = \texttt{0}
\end{itemize}

\paragraph{Interpretation.}
\begin{itemize}[leftmargin=1.5em]
  \item Increasing the collision reward weight alone did not recover alternating-impact behavior.
  \item In this sweep, larger $w_c$ actually led to an even tighter low-motion solution rather than more visible rhythmic excitation.
  \item This suggests the failure is not only that collision reward is too small; the agent may still need an intermediate shaping signal that rewards swing development before successful alternating impacts appear.
\end{itemize}

\paragraph{Status.}
\begin{itemize}[leftmargin=1.5em]
  \item This issue is not fixed yet.
  \item The next stage should continue exploring different reward mechanisms, especially intermediate shaping terms that encourage swing growth before alternating collisions are achieved.
\end{itemize}

\end{document}
